\documentclass{ximera}

\input{../../../preamble.tex}


\definecolor{fvdnavy}{HTML}{071D43}
\definecolor{fvdcyan}{HTML}{20BFC6}
\definecolor{fvdcyanlight}{HTML}{EFFCFB}
\definecolor{fvdmagenta}{HTML}{F10D69}
\definecolor{fvdmagentalight}{HTML}{FFF1F7}
\definecolor{fvdtext}{HTML}{163044}





%=================================================
% Pedagogische omgevingen
% PDF: tcolorbox
% HTML: learning-box
%=================================================

\newenvironment{voorkennis}
{%
    \ifdefined\HCode
        \HCode{%
<section class="learning-box learning-box--prior">
<div class="learning-box__header">
<span class="learning-box__icon" aria-hidden="true"></span>
<span class="learning-box__title">Voorkennis</span>
</div>
<div class="learning-box__content">
        }%
        \begin{itemize}
    \else
       \begin{tcolorbox}[
    enhanced,
    breakable,
    colback=fvdcyanlight,
    colframe=fvdcyan,
    coltext=fvdtext,
    boxrule=0.5pt,
    leftrule=4pt,
    arc=4mm,
    outer arc=4mm,
    left=7mm,
    right=6mm,
    top=5mm,
    bottom=5mm,
    before skip=6mm,
    after skip=6mm,
    title=Voorkennis,
    coltitle=fvdcyan,
    fonttitle=\bfseries\Large,
    attach title to upper={\par\smallskip},
    boxed title style={
        empty,
        left=0mm,
        right=0mm,
        top=0mm,
        bottom=0mm
    },
    overlay={
        \draw[
            fvdcyan,
            line width=0.5pt
        ]
        ([xshift=42mm,yshift=-13mm]frame.north west)
        --
        ([xshift=-7mm,yshift=-13mm]frame.north east);
    }
]
        \begin{itemize}
    \fi
}
{%
    \end{itemize}
    \ifdefined\HCode
        \HCode{%
</div>
</section>
        }%
    \else
        \end{tcolorbox}
    \fi
}


\newenvironment{leerdoelen}
{%
    \ifdefined\HCode
        \HCode{%
<section class="learning-box learning-box--goals">
<div class="learning-box__header">
<span class="learning-box__icon" aria-hidden="true"></span>
<span class="learning-box__title">Leerdoelen</span>
</div>
<div class="learning-box__content">
        }%
        \begin{itemize}
    \else
\begin{tcolorbox}[
    enhanced,
    breakable,
    colback=fvdmagentalight,
    colframe=fvdmagenta,
    coltext=fvdtext,
    boxrule=0.5pt,
    leftrule=4pt,
    arc=4mm,
    outer arc=4mm,
    left=7mm,
    right=6mm,
    top=5mm,
    bottom=5mm,
    before skip=6mm,
    after skip=6mm,
    title=Leerdoelen,
    coltitle=fvdmagenta,
    fonttitle=\bfseries\Large,
    attach title to upper={\par\smallskip},
    boxed title style={
        empty,
        left=0mm,
        right=0mm,
        top=0mm,
        bottom=0mm
    },
    overlay={
        \draw[
            fvdmagenta,
            line width=0.5pt
        ]
        ([xshift=42mm,yshift=-13mm]frame.north west)
        --
        ([xshift=-7mm,yshift=-13mm]frame.north east);
    }
]
        \begin{itemize}
    \fi
}
{%
    \end{itemize}
    \ifdefined\HCode
        \HCode{%
</div>
</section>
        }%
    \else
        \end{tcolorbox}
    \fi
}


\addPrintStyle{../../..}

\title{Elektrische kracht: aantrekken en afstoten}
\author{Ingmar Herreman}

\begin{abstract}

  Een ballon die na het wrijven aan een muur blijft hangen, haren die
  rechtop gaan staan, een kleine schok bij het aanraken van een metalen
  deurklink en een bliksemflits tijdens een onweersbui lijken op het
  eerste gezicht heel verschillende verschijnselen.

  Toch hebben ze iets gemeen: elektrische lading.

  In dit hoofdstuk onderzoeken we eerst hoe elektrische lading zich
  manifesteert. We bekijken hoe positieve en negatieve lading
  samenhangen met de atomaire structuur van materie en onderzoeken hoe
  lading kan worden overgedragen of herverdeeld.

  Daarbij maken we onderscheid tussen geleiders en isolatoren en
  bestuderen we elektrostatische influentie en aarding.

  Vervolgens onderzoeken we de krachtwerking tussen elektrische
  ladingen. Met de wet van Coulomb kunnen we de elektrische kracht
  kwantitatief beschrijven. We analyseren hoe deze kracht afhangt van
  de grootte van de ladingen en van hun onderlinge afstand.

  De wiskundige vorm van de wet van Coulomb zal opvallend vertrouwd
  lijken. In het vorige hoofdstuk vonden we voor gravitatie immers ook
  een inverse-kwadratenwet.

  Ten slotte onderzoeken we situaties waarin meerdere ladingen tegelijk
  een kracht uitoefenen. Omdat elektrische kracht een vector is, moeten
  we daarbij krachten vectorieel combineren.

  Zo maken we opnieuw de stap van waarnemen naar verklaren en van
  verklaren naar voorspellen.

\end{abstract}

\begin{document}

% FVD_CHAPTER_DATA_START
% Automatisch gegenereerd uit index.tex.
% Niet handmatig aanpassen.
\fvdchapterdata{2}{Krachten op afstand - Van gravitatie tot elektrische velden}{4}
% FVD_CHAPTER_DATA_END





\maketitle


% ============================================================
% VOORKENNIS
% ============================================================

\begin{voorkennis}
  \item Je kent de bouw van een atoom met protonen, neutronen en
        elektronen.
  \item Je kan krachten als vectoren voorstellen.
  \item Je kent de drie wetten van Newton.
  \item Je kan meerdere krachten vectorieel combineren.
  \item Je kan werken met machten van tien en wetenschappelijke notatie.
  \item Je kan SI-eenheden en voorvoegsels correct gebruiken.
  \item Je kent het begrip omgekeerd kwadratisch verband uit het
        hoofdstuk over gravitatie.
\end{voorkennis}


% ============================================================
% LEERDOELEN
% ============================================================

\begin{leerdoelen}
  \item Je kan positieve, negatieve en neutrale voorwerpen
        microscopisch verklaren.
  \item Je kan elektrische lading correct voorstellen met het symbool
        $q$ of $Q$ en de eenheid coulomb.
  \item Je kan de elementaire lading gebruiken om het aantal
        overgedragen elektronen te bepalen.
  \item Je kan het behoud van elektrische lading toepassen.
  \item Je kan verklaren hoe voorwerpen elektrisch geladen kunnen
        worden.
  \item Je kan geleiders en isolatoren onderscheiden vanuit hun
        atomaire structuur.
  \item Je kan elektrostatische influentie verklaren.
  \item Je kan aantrekking en afstoting tussen geladen voorwerpen
        analyseren.
  \item Je kan de wet van Coulomb interpreteren en toepassen.
  \item Je kan de elektrische kracht vectorieel voorstellen.
  \item Je kan Newton III toepassen op elektrische wisselwerkingen.
  \item Je kan analyseren hoe de Coulombkracht verandert wanneer
        ladingen of afstand veranderen.
  \item Je kan elektrische krachten van meerdere puntladingen
        vectorieel combineren.
  \item Je kan overeenkomsten en verschillen tussen gravitatiekracht
        en elektrische kracht uitleggen.
  \item Je kan elektrostatische verschijnselen in technologische en
        alledaagse situaties verklaren.
\end{leerdoelen}


% ============================================================
% 1. ELEKTRICITEIT ZONDER STROOM
% ============================================================

\FVDsection{Elektriciteit zonder elektrische stroom}

Elektriciteit associëren we vaak met stopcontacten, batterijen,
elektrische schakelingen en toestellen.

Maar elektrische verschijnselen kunnen ook optreden zonder dat er een
blijvende elektrische stroom door een schakeling loopt.

Denk bijvoorbeeld aan:

\begin{itemize}
  \item haren die rechtop gaan staan wanneer je een trui uittrekt;
  \item een plastic kam die papiersnippers aantrekt;
  \item een ballon die na het wrijven aan een muur blijft hangen;
  \item een kleine vonk wanneer je een metalen deurklink aanraakt;
  \item bliksem tijdens een onweersbui.
\end{itemize}

Deze verschijnselen noemen we \textbf{elektrostatische verschijnselen}.

Ze hebben te maken met elektrische ladingen en met de krachten die
elektrische ladingen op elkaar uitoefenen.


% ============================================================
% 2. EXPERIMENT
% ============================================================

\FVDsection{Een eerste experiment}

Neem een kunststofstaaf en breng die in de buurt van enkele kleine
papiersnippers.

Zonder verdere behandeling gebeurt er meestal weinig.

Wrijf de staaf vervolgens met een geschikte doek en breng ze opnieuw
in de buurt van de papiersnippers.

De snippers worden nu aangetrokken.

Er is blijkbaar iets veranderd aan de staaf.


\begin{opmerking}

We zeggen dat de staaf door het wrijven in een
\textbf{elektrisch geladen toestand} is gebracht.

Het wrijven heeft geen elektrische lading uit het niets gemaakt.
Het heeft reeds aanwezige elektrische ladingen verplaatst.

\end{opmerking}


Dat onderscheid zal verderop belangrijk worden.


% ============================================================
% 3. TWEE SOORTEN LADING
% ============================================================

\FVDsection{Twee soorten elektrische lading}

We kunnen het experiment uitbreiden.

Hang een geladen kunststofstaaf vrij op en breng een tweede, op
dezelfde manier geladen staaf in de buurt.

We stellen vast dat beide staven elkaar afstoten.

Gebruiken we twee verschillende materialen die op geschikte wijze
geladen werden, dan kunnen de staven elkaar juist aantrekken.

Uit zulke experimenten blijkt dat er twee soorten elektrische lading
bestaan.

We noemen ze:

\[
\boxed{\text{positieve lading}}
\]

en

\[
\boxed{\text{negatieve lading}}.
\]


Voor de krachtwerking geldt:

\[
\boxed{
\text{gelijknamige ladingen stoten elkaar af}
}
\]

en

\[
\boxed{
\text{tegengestelde ladingen trekken elkaar aan}.
}
\]


Schematisch:

\[
+\qquad+\qquad\Rightarrow\qquad\text{afstoting}
\]

\[
-\qquad-\qquad\Rightarrow\qquad\text{afstoting}
\]

\[
+\qquad-\qquad\Rightarrow\qquad\text{aantrekking}.
\]


Hier zien we meteen een fundamenteel verschil met gravitatie.

Bij gravitatie tussen gewone massa's vonden we alleen aantrekking.

Bij elektrische interacties zijn zowel aantrekking als afstoting
mogelijk.


% ============================================================
% 4. WAT IS ELEKTRISCHE LADING?
% ============================================================

\FVDsection{Elektrische lading}

Elektrische lading is een fundamentele eigenschap van materie.

We stellen elektrische lading meestal voor door:

\[
q
\qquad\text{of}\qquad
Q.
\]

De SI-eenheid van elektrische lading is de \textbf{coulomb}:

\[
\boxed{
[q]=\mathrm C.
}
\]


Eén coulomb is voor elektrostatische verschijnselen een grote
hoeveelheid lading. Daarom gebruiken we vaak kleinere eenheden:

\[
1\,\mathrm{mC}=10^{-3}\,\mathrm C,
\]

\[
1\,\mathrm{\mu C}=10^{-6}\,\mathrm C,
\]

\[
1\,\mathrm{nC}=10^{-9}\,\mathrm C,
\]

\[
1\,\mathrm{pC}=10^{-12}\,\mathrm C.
\]


\begin{oefening}[Voorvoegsels]{\oefB\ESS}

Zet om naar coulomb.

\begin{enumerate}
  \item $4{,}5\,\mathrm{\mu C}$
  \item $28\,\mathrm{nC}$
  \item $0{,}75\,\mathrm{mC}$
  \item $350\,\mathrm{pC}$
\end{enumerate}

\end{oefening}


% ============================================================
% 5. ATOMAIRE OORSPRONG
% ============================================================

\FVDsection{Waar komt elektrische lading vandaan?}

Om elektrostatische verschijnselen te begrijpen moeten we naar de
atomaire structuur van materie kijken.

Een atoom bestaat uit:

\begin{itemize}
  \item protonen;
  \item neutronen;
  \item elektronen.
\end{itemize}

Protonen en neutronen bevinden zich in de kern.

Elektronen bevinden zich rondom de kern.


De elektrische eigenschappen zijn:

\[
\begin{array}{c|c}
\text{deeltje} & \text{elektrische lading}\\
\hline
\text{proton} & +e\\
\text{elektron} & -e\\
\text{neutron} & 0
\end{array}
\]


De grootte van de lading van één proton of elektron noemen we de
\textbf{elementaire lading}:

\[
\boxed{
e=1{,}602\times10^{-19}\,\mathrm C.
}
\]

Dus:

\[
q_p=+e
\]

en

\[
q_e=-e.
\]


% ============================================================
% 6. NEUTRAAL, POSITIEF EN NEGATIEF
% ============================================================

\FVDsection{Wanneer is een voorwerp geladen?}

Een neutraal atoom bevat evenveel protonen als elektronen.

De totale positieve lading is dan even groot als de totale negatieve
lading.

Daardoor is:

\[
\boxed{q_\mathrm{netto}=0.}
\]


Een voorwerp wordt negatief geladen wanneer het een overschot aan
elektronen heeft.

Een voorwerp wordt positief geladen wanneer het een tekort aan
elektronen heeft.


\begin{center}
\[
\boxed{
\begin{array}{c|c}
\text{toestand} & \text{microscopische betekenis}\\
\hline
\text{neutraal}
&
N_p=N_e
\\[1mm]
\text{positief geladen}
&
N_p>N_e
\\[1mm]
\text{negatief geladen}
&
N_e>N_p
\end{array}
}
\]
\end{center}


\begin{opmerking}

Bij gewone elektrostatische processen zijn het voornamelijk
\textbf{elektronen} die van het ene voorwerp naar het andere bewegen.

Wanneer een voorwerp positief geladen wordt, betekent dit dus meestal
niet dat het extra protonen heeft opgenomen.

Het heeft elektronen afgestaan.

\end{opmerking}


% ============================================================
% 7. IONEN
% ============================================================

\FVDsection{Ionen}

Wanneer een atoom elektronen verliest of opneemt, is het niet langer
elektrisch neutraal.

Zo'n geladen atoom noemen we een \textbf{ion}.

Een atoom dat één of meerdere elektronen verliest, wordt positief
geladen.

Een atoom dat één of meerdere elektronen opneemt, wordt negatief
geladen.


\begin{oefening}[Lithium]{\oefB}

Een neutraal lithiumatoom bevat drie protonen.

\begin{enumerate}
  \item Hoeveel elektronen bevat het neutrale atoom?
  \item Het atoom staat één elektron af. Wat is het teken van de
        resulterende lading?
  \item Hoe groot is die lading?
\end{enumerate}

\end{oefening}


% ============================================================
% 8. KWANTISATIE VAN LADING
% ============================================================

\FVDsection{Elektrische lading komt in pakketjes}

Omdat elektronen een vaste elementaire lading hebben, kan de netto
lading van een voorwerp niet zomaar elke willekeurige waarde aannemen.

Voor een voorwerp dat een netto aantal $n$ elementaire ladingen heeft,
geldt:

\[
\boxed{
q=ne.
}
\]

Hierbij is $n$ een geheel getal wanneer we het teken in $n$ opnemen.

Voor de grootte van de lading kunnen we schrijven:

\[
\boxed{
|q|=Ne
}
\]

met $N$ het aantal overtollige of ontbrekende elektronen.


Dit noemen we de \textbf{kwantisatie van elektrische lading}.


\begin{voorbeeld}[Hoeveel elektronen?]

Een voorwerp heeft een lading

\[
q=-3{,}20\times10^{-6}\,\mathrm C.
\]

Het minteken vertelt ons dat er een overschot aan elektronen is.

Het aantal extra elektronen is:

\[
N=\frac{|q|}{e}.
\]

Dus:

\[
N=
\frac{3{,}20\times10^{-6}}
{1{,}602\times10^{-19}}.
\]

Daaruit volgt ongeveer:

\[
\boxed{
N=2{,}00\times10^{13}.
}
\]

Een lading van enkele microcoulomb komt dus al overeen met een enorm
aantal elektronen.

\end{voorbeeld}


\begin{oefening}[Een coulomb]{\oefB\ESS}

Hoeveel elektronen zijn nodig om een netto lading met grootte

\[
1{,}00\,\mathrm C
\]

te verkrijgen?

\end{oefening}


\begin{oefening}[Mogelijke lading?]{\oefU\ESS}

Een leerling beweert dat een klein voorwerp een netto lading heeft van

\[
q=2{,}50\times10^{-19}\,\mathrm C.
\]

Onderzoek met het model van gekwantiseerde lading of deze waarde
mogelijk is.

\end{oefening}


% ============================================================
% 9. BEHOUD VAN LADING
% ============================================================

\FVDsection{Elektrische lading wordt niet gemaakt door wrijving}

Wanneer we twee materialen over elkaar wrijven, zeggen we soms
informeel dat we ``lading opwekken''.

Dat is misleidend.

De elektrische lading was al aanwezig in de materie.

Door het wrijven kunnen elektronen van het ene materiaal naar het
andere worden overgedragen.

Het ene voorwerp krijgt daardoor een elektronenoverschot.

Het andere krijgt een even groot elektronentekort.


Voor een geïsoleerd systeem geldt het behoud van elektrische lading:

\[
\boxed{
Q_\mathrm{voor}=Q_\mathrm{na}.
}
\]


\begin{opmerking}

Elektrische lading kan binnen een systeem worden verplaatst of
gescheiden.

De totale elektrische lading van een geïsoleerd systeem blijft
behouden.

\end{opmerking}


\begin{oefening}[Ladingsoverdracht]{\oefU\ESS}

Twee aanvankelijk neutrale voorwerpen worden tegen elkaar gewreven.

Na afloop heeft voorwerp A een lading

\[
q_A=-4{,}0\,\mathrm{nC}.
\]

Neem aan dat het systeem elektrisch geïsoleerd is.

\begin{enumerate}
  \item Welke lading heeft voorwerp B?
  \item Welk voorwerp heeft elektronen opgenomen?
  \item Hoeveel elektronen zijn ongeveer overgedragen?
\end{enumerate}

\end{oefening}


% ============================================================
% 10. LADEN DOOR WRIJVING
% ============================================================

\FVDsection{Laden door wrijving}

Bij contact tussen verschillende materialen kunnen elektronen worden
overgedragen.

Door wrijving wordt het contact tussen de oppervlakken sterk vergroot
en wordt ladingsscheiding vaak duidelijk waarneembaar.

Wanneer materiaal A elektronen overdraagt aan materiaal B:

\[
A\longrightarrow A^+,
\]

\[
B\longrightarrow B^-.
\]

Het is dus niet de wrijving zelf die nieuwe lading creëert.

De wrijving helpt bestaande ladingen te scheiden.


% ============================================================
% 11. LADEN DOOR CONTACT
% ============================================================

\FVDsection{Laden door contact}

Een geladen voorwerp kan ook lading overdragen door rechtstreeks
contact met een ander voorwerp.

Stel dat een negatief geladen metalen bol een neutrale metalen bol
aanraakt.

Omdat elektronen zich in metaal gemakkelijk kunnen verplaatsen, kunnen
een aantal overtollige elektronen zich over beide bollen verdelen.

Na het contact kunnen beide bollen negatief geladen zijn.

De totale lading blijft daarbij behouden.


% ============================================================
% 12. GELEIDERS EN ISOLATOREN
% ============================================================

\FVDsection{Waarom beweegt lading in sommige materialen gemakkelijker?}

Niet elk materiaal reageert op dezelfde manier op elektrische lading.

We onderscheiden twee belangrijke groepen.


\begin{definitie}[Geleider]

Een \textbf{geleider} is een materiaal waarin elektrische ladingen zich
gemakkelijk door het materiaal kunnen verplaatsen.

\end{definitie}


\begin{definitie}[Isolator]

Een \textbf{isolator} is een materiaal waarin elektrische ladingen zich
niet gemakkelijk over grote afstanden door het materiaal kunnen
verplaatsen.

\end{definitie}


Metalen zijn goede geleiders.

Voorbeelden van isolatoren zijn onder andere glas, droge lucht en veel
kunststoffen.


% ============================================================
% 13. ATOMAIRE VERKLARING GELEIDER
% ============================================================

\FVDsection{Een microscopisch model van een metalen geleider}

In een metaal zijn sommige buitenste elektronen relatief zwak aan één
specifiek atoom gebonden.

We kunnen een metaal vereenvoudigd voorstellen als een rooster van
positieve ionen waarin elektronen zich relatief vrij kunnen bewegen.

Deze mobiele elektronen noemen we vaak \textbf{vrije elektronen} of
\textbf{geleidingselektronen}.

Wanneer ergens op een metalen voorwerp een overschot aan elektronen
wordt aangebracht, kunnen deze elektronen zich over het voorwerp
herverdelen.

Bij een isolator zijn de relevante elektronen veel sterker lokaal
gebonden.

Een aangebrachte lading blijft daardoor veel meer in de buurt van de
plaats waar ze werd aangebracht.


\begin{oefening}[Geleider of isolator?]{\oefB\ESS}

Leg voor elk van de volgende uitspraken uit of ze past bij een
geleider of een isolator.

\begin{enumerate}
  \item Overtollige elektronen kunnen zich gemakkelijk herverdelen.
  \item Een aangebrachte lading blijft grotendeels lokaal.
  \item Het materiaal bevat mobiele ladingsdragers.
  \item Het materiaal wordt gebruikt om een elektrische draad van de
        omgeving af te schermen.
\end{enumerate}

\end{oefening}


% ============================================================
% 14. INFLUENTIE
% ============================================================

\FVDsection{Ladingen beïnvloeden elkaar zonder contact}

Breng een positief geladen voorwerp in de buurt van een neutrale
metalen geleider.

De twee voorwerpen raken elkaar niet.

Toch gebeurt er iets in de geleider.

De vrije elektronen worden door de positieve externe lading
aangetrokken.

Ze verplaatsen zich naar de zijde die het dichtst bij de positieve
lading ligt.

Aan die kant ontstaat een elektronenoverschot.

Aan de andere kant ontstaat een elektronentekort.

De geleider als geheel kan nog steeds neutraal zijn, maar de ladingen
zijn niet langer gelijkmatig verdeeld.

Dit noemen we \textbf{elektrostatische influentie}.


\begin{definitie}[Elektrostatische influentie]

Elektrostatische influentie is de herverdeling van elektrische lading
in een voorwerp onder invloed van een nabijgelegen elektrische lading,
zonder rechtstreeks contact met die externe lading.

\end{definitie}


% ============================================================
% 15. AANTREKKING VAN EEN NEUTRAAL VOORWERP
% ============================================================

\FVDsection{Kan een geladen voorwerp een neutraal voorwerp aantrekken?}

Ja.

Dat is een belangrijke nuance.

``Neutraal'' betekent:

\[
Q_\mathrm{netto}=0.
\]

Het betekent niet dat het voorwerp geen positieve en negatieve
ladingen bevat.

Wanneer een geladen voorwerp in de buurt komt, kunnen de ladingen in
het neutrale voorwerp enigszins verschuiven of polariseren.

De tegengestelde lading bevindt zich gemiddeld dichter bij de externe
lading dan de gelijknamige lading.

Omdat elektrische kracht sterk afhangt van de afstand, kan daardoor
een netto aantrekkingskracht ontstaan.


\begin{opmerking}

Uit het feit dat twee voorwerpen elkaar aantrekken, mag je dus niet
automatisch besluiten dat ze tegengesteld geladen zijn.

Eén van beide kan elektrisch neutraal zijn.

Afstoting is in dat opzicht een sterkere aanwijzing voor de tekens van
de ladingen.

\end{opmerking}


\begin{oefening}[Een geladen kam]{\oefU\ESS}

Een negatief geladen plastic kam trekt neutrale papiersnippers aan.

Een leerling zegt:

``De papiersnippers moeten dus positief geladen zijn.''

Analyseer deze uitspraak en geef een betere verklaring.

\end{oefening}


% ============================================================
% 16. AARDING
% ============================================================

\FVDsection{Aarding}

De aarde is zo groot dat ze in veel elektrostatische situaties als een
enorm reservoir van elektrische lading kan worden beschouwd.

Wanneer een geleider elektrisch met de aarde verbonden wordt, kunnen
elektronen van of naar de aarde bewegen.

Dit noemen we \textbf{aarding}.

Aarding kan daardoor worden gebruikt om ongewenste ladingsophoping te
verminderen of om een voorwerp via influentie te laden.


% ============================================================
% 17. LADEN DOOR INFLUENTIE
% ============================================================

\FVDsection{Laden zonder aanraken}

Een geleider kan zelfs een netto lading krijgen zonder rechtstreeks
contact met het geladen voorwerp dat de ladingsscheiding veroorzaakt.

Een mogelijke procedure is:

\begin{enumerate}
  \item breng een geladen voorwerp dicht bij een neutrale geleider;
  \item de vrije elektronen in de geleider herverdelen zich;
  \item verbind de geleider tijdelijk met de aarde;
  \item elektronen kunnen naar of vanuit de aarde bewegen;
  \item verbreek eerst de aardverbinding;
  \item verwijder daarna het externe geladen voorwerp.
\end{enumerate}

De geleider kan dan netto geladen achterblijven.

Dit is laden door \textbf{influentialading}.


% ============================================================
% 18. STEM: STATISCHE ELEKTRICITEIT
% ============================================================

\FVDsection{STEM-toepassing: wanneer wordt statische elektriciteit gevaarlijk?}

Een kleine elektrostatische vonk is thuis meestal alleen vervelend.

In een industriële omgeving kan dezelfde fysica echter ernstige
gevolgen hebben.

Bij het verpompen van brandstoffen, het transporteren van poeders of
het bewegen van kunststofbanden kunnen door contact en wrijving
ladingen gescheiden worden.

Wanneer die lading niet voldoende kan wegstromen, kan een groot
potentiaalverschil ontstaan.

Een ontlading kan dan een vonk veroorzaken.

In een omgeving met brandbare dampen of fijn stof kan zo'n vonk een
ontsteking veroorzaken.

Daarom worden in industriële installaties onder andere:

\begin{itemize}
  \item geleidende verbindingen gebruikt;
  \item metalen onderdelen met elkaar verbonden;
  \item installaties geaard;
  \item ongewenste ladingsophopingen vermeden.
\end{itemize}


Dit is een mooi voorbeeld van hoe een fundamenteel fysisch concept
rechtstreeks samenhangt met veiligheid en engineering.


% ============================================================
% 19. VAN KWALITATIEF NAAR KWANTITATIEF
% ============================================================

\FVDsection{Hoe groot is de elektrische kracht?}

Tot nu toe konden we kwalitatief voorspellen of twee ladingen elkaar
aantrekken of afstoten.

Maar fysica wil meer.

We willen ook voorspellen:

\begin{center}
\textit{hoe groot is de kracht?}
\end{center}

Daarvoor gebruiken we de \textbf{wet van Coulomb}.


% ============================================================
% 20. PUNTLADINGMODEL
% ============================================================

\FVDsection{Het puntladingmodel}

Net zoals we bij gravitatie massa's soms als puntmassa's modelleerden,
kunnen we geladen voorwerpen soms als \textbf{puntladingen}
beschouwen.

Een puntlading is een ideaal model waarbij we de afmetingen van het
geladen voorwerp verwaarlozen ten opzichte van de afstanden in het
probleem.

Dat model is niet altijd geldig.

Voor uitgebreide geladen voorwerpen met een complexe ladingsverdeling
kan de situatie veel ingewikkelder worden.


% ============================================================
% 21. WET VAN COULOMB
% ============================================================

\FVDsection{De wet van Coulomb}

Beschouw twee puntladingen $Q_1$ en $Q_2$ op onderlinge afstand $r$.

De grootte van de elektrische kracht tussen beide ladingen wordt
gegeven door:


\begin{definitie}[Wet van Coulomb]

\[
\boxed{
F_e
=
k\frac{|Q_1Q_2|}{r^2}.
}
\]

Hierin is:

\[
F_e=\text{grootte van de elektrische kracht},
\]

\[
Q_1,Q_2=\text{elektrische ladingen},
\]

\[
r=\text{afstand tussen de puntladingen},
\]

en in vacuüm:

\[
\boxed{
k\approx8{,}99\times10^9\,
\frac{\mathrm{N\,m^2}}{\mathrm{C^2}}.
}
\]

\end{definitie}


De absolute waarde in

\[
|Q_1Q_2|
\]

zorgt ervoor dat de formule een positieve \textbf{grootte} voor de
kracht oplevert.

Het teken van de ladingen gebruiken we afzonderlijk om te bepalen of
de kracht aantrekkend of afstotend is.


% ============================================================
% 22. DE KRACHT ALS VECTOR
% ============================================================

\FVDsection{De elektrische kracht is een vector}

De wet van Coulomb geeft de grootte van de kracht.

Maar kracht is een vector.

We moeten dus ook richting en zin bepalen.

Voor twee puntladingen ligt de richting van de elektrische kracht langs
de verbindingslijn tussen beide ladingen.


Bij gelijknamige ladingen:

\[
+\quad\longleftarrow
\qquad
\longrightarrow\quad+
\]

of:

\[
-\quad\longleftarrow
\qquad
\longrightarrow\quad-
\]

De ladingen stoten elkaar af.


Bij tegengestelde ladingen:

\[
+\quad\longrightarrow
\qquad
\longleftarrow\quad-
\]

De ladingen trekken elkaar aan.


Gebruik daarom bij elk Coulombprobleem twee afzonderlijke vragen:

\[
\boxed{\text{Hoe groot is de kracht?}}
\]

en:

\[
\boxed{\text{In welke richting en zin werkt de kracht?}}
\]


% ============================================================
% 23. NEWTON III
% ============================================================

\FVDsection{Welke lading oefent de grootste kracht uit?}

Stel:

\[
Q_1=+1\,\mathrm{\mu C}
\]

en:

\[
Q_2=-10\,\mathrm{\mu C}.
\]

Omdat $|Q_2|$ tienmaal groter is dan $|Q_1|$, zou je kunnen denken dat
$Q_2$ een grotere kracht op $Q_1$ uitoefent.

Dat is fout.

Volgens Newton III:

\[
\boxed{
\vec F_{1\rightarrow2}
=
-\vec F_{2\rightarrow1}.
}
\]

De krachten hebben:

\begin{itemize}
  \item dezelfde grootte;
  \item dezelfde richting;
  \item tegengestelde zin;
  \item een verschillend aangrijpingspunt.
\end{itemize}


De wet van Coulomb bevestigt dat:

\[
F_e
=
k\frac{|Q_1Q_2|}{r^2}.
\]

Het product

\[
|Q_1Q_2|
\]

is voor beide krachten hetzelfde.


\begin{oefening}[Newton III]{\oefU\ESS}

Twee puntladingen zijn:

\[
Q_1=+2{,}0\,\mathrm{\mu C}
\]

en:

\[
Q_2=-8{,}0\,\mathrm{\mu C}.
\]

Een leerling zegt:

``De lading van $-8{,}0\,\mathrm{\mu C}$ trekt viermaal harder aan de
andere lading.''

Analyseer de uitspraak.

\end{oefening}


% ============================================================
% 24. EENHEDENCONTROLE
% ============================================================

\FVDsection{Eenheden controleren}

Uit

\[
F_e=k\frac{|Q_1Q_2|}{r^2}
\]

volgt:

\[
[F_e]
=
\frac{\mathrm{N\,m^2}}{\mathrm{C^2}}
\frac{\mathrm{C^2}}{\mathrm{m^2}}.
\]

Dus:

\[
\boxed{
[F_e]=\mathrm N.
}
\]

Ook hier vormen eenheden een krachtig controle-instrument.


% ============================================================
% 25. VOORBEELD COULOMB
% ============================================================

\FVDsection{Een eerste Coulombberekening}

\begin{voorbeeld}[Twee tegengestelde ladingen]

Twee puntladingen zijn:

\[
Q_1=+2{,}0\,\mathrm{\mu C}
\]

en:

\[
Q_2=-3{,}0\,\mathrm{\mu C}.
\]

De afstand tussen beide bedraagt:

\[
r=0{,}20\,\mathrm m.
\]

We zoeken eerst de grootte van de elektrische kracht:

\[
F_e
=
k\frac{|Q_1Q_2|}{r^2}.
\]

Zet de ladingen om naar coulomb:

\[
Q_1=2{,}0\times10^{-6}\,\mathrm C,
\]

\[
Q_2=-3{,}0\times10^{-6}\,\mathrm C.
\]

Invullen geeft:

\[
F_e
=
8{,}99\times10^9
\frac{
(2{,}0\times10^{-6})
(3{,}0\times10^{-6})
}{
(0{,}20)^2
}.
\]

Dus:

\[
\boxed{
F_e\approx1{,}35\,\mathrm N.
}
\]

Omdat de ladingen tegengesteld zijn, is de kracht
\textbf{aantrekkend}.

De krachtvector op elke lading wijst dus naar de andere lading.

\end{voorbeeld}


% ============================================================
% 26. SYSTEMATISCHE AANPAK
% ============================================================

\FVDsection{Een Coulombprobleem systematisch oplossen}

Bij eenvoudige oefeningen lijkt rechtstreeks invullen aantrekkelijk.

Bij meerdere ladingen werkt dat niet meer betrouwbaar.

Gebruik daarom steeds dezelfde aanpak:

\[
\boxed{
\begin{array}{c}
\text{maak een schets}
\\
\downarrow
\\
\text{noteer gegeven en gevraagd}
\\
\downarrow
\\
\text{zet alle grootheden om naar SI-eenheden}
\\
\downarrow
\\
\text{bepaal aantrekking of afstoting}
\\
\downarrow
\\
\text{teken de krachtvector(en)}
\\
\downarrow
\\
\text{bereken de grootte(n)}
\\
\downarrow
\\
\text{combineer vectorieel indien nodig}
\\
\downarrow
\\
\text{controleer eenheid en grootteorde}
\\
\downarrow
\\
\text{interpreteer het antwoord}
\end{array}
}
\]


Deze volgorde voorkomt een veel voorkomende fout:

\begin{center}
eerst rekenen en pas daarna proberen te bedenken in welke richting de
kracht werkt.
\end{center}


Bij vectorproblemen moet de fysische analyse juist \textbf{voor} de
berekening gebeuren.


% ============================================================
% 27. INVLOED VAN DE LADINGEN
% ============================================================

\FVDsection{Hoe beïnvloedt de grootte van de lading de kracht?}

Uit:

\[
F_e
=
k\frac{|Q_1Q_2|}{r^2}
\]

volgt bij constante afstand:

\[
\boxed{
F_e\propto |Q_1Q_2|.
}
\]

Wanneer één lading verdubbelt:

\[
Q_1\rightarrow2Q_1,
\]

dan:

\[
F_e\rightarrow2F_e.
\]

Wanneer beide ladingen verdubbelen:

\[
Q_1\rightarrow2Q_1,
\qquad
Q_2\rightarrow2Q_2,
\]

dan:

\[
F_e\rightarrow4F_e.
\]


\begin{oefening}[Redeneren met verhoudingen]{\oefB\ESS}

Twee puntladingen oefenen een kracht met grootte $F$ op elkaar uit.

Bepaal de nieuwe kracht wanneer:

\begin{enumerate}
  \item $|Q_1|$ driemaal zo groot wordt;
  \item beide ladingen tweemaal zo groot worden;
  \item $|Q_1|$ halveert en $|Q_2|$ verdubbelt;
  \item beide ladingen gehalveerd worden.
\end{enumerate}

Bereken niet opnieuw met de constante $k$.

\end{oefening}


% ============================================================
% 28. INVERSE-KWADRATENWET
% ============================================================

\FVDsection{De invloed van de afstand}

Uit de wet van Coulomb volgt:

\[
\boxed{
F_e\propto\frac{1}{r^2}.
}
\]

Ook de elektrische kracht volgt dus een
\textbf{inverse-kwadratenwet}.

Als:

\[
r\rightarrow2r,
\]

dan:

\[
F_e\rightarrow\frac14F_e.
\]

Als:

\[
r\rightarrow3r,
\]

dan:

\[
F_e\rightarrow\frac19F_e.
\]

Als:

\[
r\rightarrow\frac r2,
\]

dan:

\[
F_e\rightarrow4F_e.
\]


\begin{center}
\[
\boxed{
\begin{array}{c|c}
\text{afstand} & \text{elektrische kracht}\\
\hline
r & F_e\\
2r & F_e/4\\
3r & F_e/9\\
4r & F_e/16\\
r/2 & 4F_e
\end{array}
}
\]
\end{center}


\begin{oefening}[Afstand]{\oefB\ESS}

Twee ladingen oefenen een elektrische kracht $F$ op elkaar uit.

Bepaal de nieuwe kracht wanneer de afstand:

\begin{enumerate}
  \item tweemaal zo groot wordt;
  \item vijfmaal zo groot wordt;
  \item driemaal zo klein wordt.
\end{enumerate}

\end{oefening}


% ============================================================
% 29. FUNCTIE EN GRAFIEK
% ============================================================

\FVDsection{De Coulombkracht als functie}

Voor twee vaste ladingen kunnen we schrijven:

\[
F_e(r)=\frac{C}{r^2},
\]

waarbij:

\[
C=k|Q_1Q_2|.
\]

De elektrische kracht neemt dus niet lineair af met de afstand.

Dat is belangrijk.

Een verdubbeling van de afstand betekent niet een halvering van de
kracht, maar:

\[
\boxed{
F_e(2r)=\frac14F_e(r).
}
\]


Wiskundig geldt:

\[
\lim_{r\rightarrow\infty}F_e(r)=0.
\]

De kracht wordt dus steeds kleiner wanneer de afstand toeneemt.


\begin{oefening}[Een verkeerde grafiek]{\oefU\ESS}

Een leerling tekent de Coulombkracht als functie van de afstand als
een dalende rechte.

Analyseer waarom die grafiek niet overeenkomt met:

\[
F_e(r)=\frac{C}{r^2}.
\]

Geef minstens twee eigenschappen waaraan een correcte grafiek moet
voldoen.

\end{oefening}


% ============================================================
% 30. COMBINATIE VAN VERANDERINGEN
% ============================================================

\FVDsection{Meerdere veranderingen tegelijk}

Bij analyseren moeten we verschillende effecten kunnen combineren.

Stel bijvoorbeeld:

\[
Q_1'=2Q_1,
\]

\[
Q_2'=3Q_2,
\]

en:

\[
r'=2r.
\]

Dan:

\[
F_e'
=
k\frac{|(2Q_1)(3Q_2)|}{(2r)^2}.
\]

Dus:

\[
F_e'
=
\frac{6}{4}
k\frac{|Q_1Q_2|}{r^2}.
\]

Daarom:

\[
\boxed{
F_e'=\frac32F_e.
}
\]


\begin{oefening}[Analyse zonder getallen]{\oefU\ESS}

Twee ladingen oefenen aanvankelijk een kracht $F$ op elkaar uit.

Daarna wordt:

\[
Q_1'=4Q_1,
\qquad
Q_2'=\frac12Q_2,
\qquad
r'=3r.
\]

Bepaal:

\[
\frac{F_e'}{F_e}.
\]

Werk met verhoudingen en gebruik geen numerieke waarde voor $k$.

\end{oefening}


% ============================================================
% 31. OMGEKEERDE PROBLEMEN
% ============================================================

\FVDsection{Niet altijd is de kracht gevraagd}

De wet van Coulomb kan ook gebruikt worden wanneer een lading of een
afstand onbekend is.

Uit:

\[
F_e=k\frac{|Q_1Q_2|}{r^2}
\]

kunnen we bijvoorbeeld $r$ bepalen:

\[
r^2=
k\frac{|Q_1Q_2|}{F_e}.
\]

Dus:

\[
\boxed{
r=
\sqrt{
k\frac{|Q_1Q_2|}{F_e}
}.
}
\]


Wanneer beide ladingen dezelfde grootte $Q$ hebben:

\[
F_e=k\frac{Q^2}{r^2},
\]

zodat:

\[
\boxed{
Q=r\sqrt{\frac{F_e}{k}}.
}
\]


\begin{oefening}[Twee gelijke ladingen]{\oefB\ESS}

Twee gelijke ladingen bevinden zich op een afstand van

\[
10{,}0\,\mathrm{cm}
\]

en stoten elkaar af met een kracht van:

\[
0{,}100\,\mathrm N.
\]

Bereken de grootte van elke lading.

\end{oefening}


% ============================================================
% 32. PROTON EN ELEKTRON
% ============================================================

\FVDsection{Elektrische kracht op atomaire schaal}

Beschouw een proton en een elektron.

Hun ladingen zijn:

\[
q_p=+e
\]

en:

\[
q_e=-e.
\]

De grootte van de elektrische kracht tussen beide is:

\[
F_e
=
k\frac{e^2}{r^2}.
\]

Op atomaire afstanden kan deze kracht, ondanks de zeer kleine
ladingen, fysisch zeer belangrijk zijn.

Dat komt doordat de afstand $r$ eveneens extreem klein is.


\begin{oefening}[Proton en elektron]{\oefB}

Een proton en elektron bevinden zich op een onderlinge afstand:

\[
r=5{,}0\times10^{-11}\,\mathrm m.
\]

Bereken de grootte van de elektrische kracht tussen beide.

Bepaal ook of de kracht aantrekkend of afstotend is.

\end{oefening}


% ============================================================
% 33. VERGELIJKING MET GRAVITATIE
% ============================================================

\FVDsection{Coulomb en Newton: een opvallende overeenkomst}

In het vorige hoofdstuk vonden we:

\[
\boxed{
F_g
=
G\frac{m_1m_2}{r^2}.
}
\]

Nu vinden we:

\[
\boxed{
F_e
=
k\frac{|Q_1Q_2|}{r^2}.
}
\]

De wiskundige structuur is opvallend gelijkaardig.


\begin{center}
\[
\boxed{
\begin{array}{c|c}
\text{gravitatie} & \text{elektrische interactie}\\
\hline
m & Q\\[1mm]
G & k\\[1mm]
F_g=G\dfrac{m_1m_2}{r^2}
&
F_e=k\dfrac{|Q_1Q_2|}{r^2}\\[3mm]
\text{inverse-kwadratenwet}
&
\text{inverse-kwadratenwet}\\[1mm]
\text{aantrekkend}
&
\text{aantrekkend of afstotend}
\end{array}
}
\]
\end{center}


Er is echter een fundamenteel verschil.

Voor gewone materie kennen we positieve massa, waardoor gravitatie
aantrekkend is.

Elektrische lading bestaat in twee tekens.

Daarom kan de elektrische wisselwerking aantrekken of afstoten.


% ============================================================
% 34. HOEVEEL STERKER?
% ============================================================

\FVDsection{Elektriciteit versus gravitatie op deeltjesniveau}

We kunnen voor een proton en een elektron zowel de elektrische als de
gravitationele kracht bekijken.

Elektrisch:

\[
F_e
=
k\frac{e^2}{r^2}.
\]

Gravitationeel:

\[
F_g
=
G\frac{m_pm_e}{r^2}.
\]

Neem de verhouding:

\[
\frac{F_e}{F_g}
=
\frac{
k\dfrac{e^2}{r^2}
}{
G\dfrac{m_pm_e}{r^2}
}.
\]

De afstand valt weg:

\[
\boxed{
\frac{F_e}{F_g}
=
\frac{ke^2}{Gm_pm_e}.
}
\]

De verhouding is dus onafhankelijk van de onderlinge afstand.


\begin{oefening}[Twee fundamentele krachten]{\oefG}

Gebruik:

\[
k=8{,}99\times10^9\,
\frac{\mathrm{N\,m^2}}{\mathrm{C^2}},
\]

\[
e=1{,}602\times10^{-19}\,\mathrm C,
\]

\[
G=6{,}674\times10^{-11}\,
\frac{\mathrm{N\,m^2}}{\mathrm{kg^2}},
\]

\[
m_p=1{,}673\times10^{-27}\,\mathrm{kg},
\]

\[
m_e=9{,}109\times10^{-31}\,\mathrm{kg}.
\]

Bereken:

\[
\frac{F_e}{F_g}.
\]

Interpreteer de grootteorde van je resultaat.

Leg ook uit waarom de afstand niet nodig is om deze verhouding te
berekenen.

\end{oefening}


% ============================================================
% 35. MEERDERE LADINGEN
% ============================================================

\FVDsection{Wat als er meer dan twee ladingen zijn?}

De wet van Coulomb beschrijft de kracht tussen twee puntladingen.

Maar een lading kan tegelijk krachten ondervinden van meerdere andere
ladingen.

Elke bronlading veroorzaakt afzonderlijk een elektrische kracht.

De resulterende elektrische kracht vinden we door de afzonderlijke
krachtvectoren op te tellen:

\[
\boxed{
\vec F_\mathrm{res}
=
\vec F_1+\vec F_2+\vec F_3+\cdots
}
\]

Dit noemen we het \textbf{superpositiebeginsel}.


% ============================================================
% 36. DRIE LADINGEN OP EEN RECHTE
% ============================================================

\FVDsection{Superpositie op één rechte}

Beschouw drie ladingen op één rechte:

\[
Q_1
\qquad
Q_2
\qquad
Q_3.
\]

We willen bijvoorbeeld de resulterende kracht op $Q_3$ bepalen.

Dan onderzoeken we afzonderlijk:

\[
\vec F_{1\rightarrow3}
\]

en:

\[
\vec F_{2\rightarrow3}.
\]

Daarna:

\[
\boxed{
\vec F_{\mathrm{res},3}
=
\vec F_{1\rightarrow3}
+
\vec F_{2\rightarrow3}.
}
\]


Het belangrijkste deel van de oplossing gebeurt vóór de berekening:

\begin{enumerate}
  \item bepaal voor elke interactie aantrekking of afstoting;
  \item teken elke krachtvector op de juiste lading;
  \item bepaal voor elke kracht de correcte afstand;
  \item bereken pas daarna de groottes;
  \item tel de krachten vectorieel op.
\end{enumerate}


% ============================================================
% 37. VOORBEELD DRIE LADINGEN
% ============================================================

\begin{voorbeeld}[Drie ladingen op één rechte]

Drie puntladingen liggen op één rechte.

\[
Q_1=-5{,}00\,\mathrm{\mu C},
\]

\[
Q_2=+10{,}0\,\mathrm{\mu C},
\]

\[
Q_3=-5{,}00\,\mathrm{\mu C}.
\]

De afstanden zijn:

\[
r_{12}=5{,}00\,\mathrm{cm},
\]

\[
r_{23}=2{,}50\,\mathrm{cm}.
\]

We zoeken de resulterende kracht op $Q_3$.

Eerst bepalen we:

\[
r_{13}=r_{12}+r_{23}.
\]

Dus:

\[
r_{13}=7{,}50\,\mathrm{cm}.
\]

Omdat $Q_1$ en $Q_3$ beide negatief zijn, stoot $Q_1$ de lading
$Q_3$ af.

Omdat $Q_2$ positief is en $Q_3$ negatief, trekt $Q_2$ de lading
$Q_3$ aan.

De twee krachten hebben dus tegengestelde zin.

We berekenen afzonderlijk:

\[
F_{1\rightarrow3}
=
k\frac{|Q_1Q_3|}{r_{13}^2}
\]

en:

\[
F_{2\rightarrow3}
=
k\frac{|Q_2Q_3|}{r_{23}^2}.
\]

Daarna combineren we beide krachten met hun correcte zin.

Het teken van de ladingen wordt dus niet gebruikt om een negatieve
kracht uit de formule te laten verschijnen.

De tekens worden gebruikt om de \textbf{richting en zin} van de
krachtvectoren te bepalen.

\end{voorbeeld}


% ============================================================
% 38. TWEE DIMENSIES
% ============================================================

\FVDsection{Elektrische krachten in twee dimensies}

Wanneer de ladingen niet op één rechte liggen, volstaat optellen of
aftrekken van de groottes niet.

We moeten de krachten werkelijk als vectoren behandelen.

Stel dat:

\[
\vec F_1
\]

en:

\[
\vec F_2
\]

onder een hoek op elkaar staan.

Dan geldt nog steeds:

\[
\vec F_\mathrm{res}
=
\vec F_1+\vec F_2.
\]

We kunnen de krachten ontbinden in componenten:

\[
F_x=\sum_i F_{i,x},
\]

\[
F_y=\sum_i F_{i,y}.
\]

Daaruit:

\[
\boxed{
F_\mathrm{res}
=
\sqrt{F_x^2+F_y^2}.
}
\]

De richting volgt bijvoorbeeld uit:

\[
\boxed{
\tan\theta=\frac{F_y}{F_x},
}
\]

waarbij we bij het bepalen van de hoek steeds rekening houden met het
juiste kwadrant.


% ============================================================
% 39. SYMMETRIE
% ============================================================

\FVDsection{Gebruik symmetrie voordat je rekent}

Bij sommige configuraties kunnen we veel werk besparen door eerst naar
de geometrie te kijken.

Stel dat twee gelijke ladingen symmetrisch ten opzichte van een derde
lading liggen.

Dan kunnen bepaalde componenten van de krachten elkaar opheffen,
terwijl andere componenten elkaar versterken.

Een goede fysicus begint daarom niet altijd met een formule.

Soms is de eerste vraag:

\begin{center}
\textit{Welke symmetrie zit in deze situatie?}
\end{center}


\begin{oefening}[Symmetrische configuratie]{\oefU\ESS}

Twee gelijke positieve ladingen $Q$ bevinden zich symmetrisch links en
rechts van een positieve testlading $q$.

De afstanden tot $q$ zijn gelijk.

\begin{enumerate}
  \item Teken beide krachtvectoren op $q$.
  \item Vergelijk hun groottes.
  \item Bepaal de resulterende kracht.
  \item Leg uit waarom hiervoor geen numerieke berekening nodig is.
\end{enumerate}

\end{oefening}


% ============================================================
% 40. OEFENINGENBUNDEL
% ============================================================

\FVDsection{Oefeningen}


% ------------------------------------------------------------
% BASIS
% ------------------------------------------------------------

\FVDsubsection{Basis}


\begin{oefening}[Soorten lading]{\oefB}

Bepaal voor elke combinatie of er aantrekking of afstoting optreedt.

\begin{enumerate}
  \item $Q_1>0$ en $Q_2>0$
  \item $Q_1<0$ en $Q_2<0$
  \item $Q_1>0$ en $Q_2<0$
  \item $Q_1<0$ en $Q_2>0$
\end{enumerate}

\end{oefening}


\begin{oefening}[Elektronenoverschot]{\oefB\ESS}

Een voorwerp heeft een lading:

\[
q=-8{,}01\times10^{-9}\,\mathrm C.
\]

\begin{enumerate}
  \item Heeft het voorwerp een tekort of overschot aan elektronen?
  \item Bereken hoeveel elektronen daarmee overeenkomen.
\end{enumerate}

\end{oefening}


\begin{oefening}[Elektronentekort]{\oefB\ESS}

Een metalen bol heeft een positieve lading:

\[
q=+4{,}80\,\mathrm{nC}.
\]

\begin{enumerate}
  \item Heeft de bol elektronen opgenomen of afgestaan?
  \item Bereken het aantal ontbrekende elektronen.
\end{enumerate}

\end{oefening}


\begin{oefening}[Coulombkracht]{\oefB\ESS}

Twee puntladingen:

\[
Q_1=+4{,}0\,\mathrm{\mu C}
\]

en:

\[
Q_2=-6{,}0\,\mathrm{\mu C}
\]

bevinden zich op een afstand:

\[
r=0{,}30\,\mathrm m.
\]

\begin{enumerate}
  \item Bereken de grootte van de elektrische kracht.
  \item Is de kracht aantrekkend of afstotend?
  \item Teken beide krachtvectoren.
\end{enumerate}

\end{oefening}


\begin{oefening}[Afstand verdriedubbelen]{\oefB\ESS}

Twee ladingen oefenen een elektrische kracht met grootte $F$ op elkaar
uit.

De afstand wordt verdrievoudigd terwijl de ladingen gelijk blijven.

Bepaal de nieuwe kracht in functie van $F$.

\end{oefening}


\begin{oefening}[Onbekende lading]{\oefB\ESS}

Twee identieke puntladingen bevinden zich op:

\[
r=1{,}00\,\mathrm m
\]

van elkaar en oefenen een kracht uit van:

\[
F_e=1{,}00\,\mathrm N.
\]

Bereken de grootte van elke lading.

\end{oefening}


% ------------------------------------------------------------
% VERDIEPING
% ------------------------------------------------------------

\FVDsubsection{Verdieping}


\begin{oefening}[Wat kan je besluiten?]{\oefU\ESS}

Twee kleine voorwerpen trekken elkaar aan.

Een leerling besluit:

``De ene lading is positief en de andere negatief.''

Is die conclusie noodzakelijk juist?

Bespreek minstens twee mogelijke situaties die tot aantrekking kunnen
leiden.

\end{oefening}


\begin{oefening}[Veranderde situatie]{\oefU\ESS}

Twee puntladingen oefenen een kracht $F$ op elkaar uit.

Vervolgens:

\[
Q_1'=3Q_1,
\]

\[
Q_2'=2Q_2,
\]

en:

\[
r'=\frac32r.
\]

Bepaal:

\[
\frac{F_e'}{F_e}.
\]

\end{oefening}


\begin{oefening}[Van kracht naar afstand]{\oefU\ESS}

Twee puntladingen hebben groottes:

\[
|Q_1|=2{,}0\,\mathrm{\mu C}
\]

en:

\[
|Q_2|=5{,}0\,\mathrm{\mu C}.
\]

De grootte van de kracht tussen beide is:

\[
F_e=0{,}50\,\mathrm N.
\]

Bereken de onderlinge afstand.

Kan je uit deze gegevens alleen bepalen of de kracht aantrekkend of
afstotend is? Verklaar.

\end{oefening}


\begin{oefening}[Drie ladingen op een rechte]{\oefU\ESS}

Drie puntladingen liggen op de $x$-as.

\[
Q_1=+2{,}0\,\mathrm{\mu C},
\]

\[
Q_2=-3{,}0\,\mathrm{\mu C},
\]

\[
Q_3=+4{,}0\,\mathrm{\mu C}.
\]

Hun posities zijn:

\[
x_1=0,
\qquad
x_2=0{,}20\,\mathrm m,
\qquad
x_3=0{,}50\,\mathrm m.
\]

Bepaal de resulterende elektrische kracht op $Q_2$.

Werk systematisch:

\begin{enumerate}
  \item maak een schets;
  \item teken de afzonderlijke krachtvectoren;
  \item bepaal de correcte afstanden;
  \item bereken de groottes;
  \item bepaal de resulterende kracht met richting en zin.
\end{enumerate}

\end{oefening}


\begin{oefening}[Fout zoeken]{\oefU\ESS}

Een leerling berekent:

\[
F_e
=
8{,}99\times10^9
\frac{
(-2{,}0\times10^{-6})
(+5{,}0\times10^{-6})
}{
(0{,}10)^2
}
\]

en vindt:

\[
F_e=-8{,}99\,\mathrm N.
\]

Daarna schrijft de leerling:

``De elektrische kracht heeft dus een negatieve grootte.''

Analyseer de redenering.

Geef aan:

\begin{enumerate}
  \item wat de fout is;
  \item hoe de grootte correct wordt berekend;
  \item welke informatie de tekens van de ladingen wel geven.
\end{enumerate}

\end{oefening}


% ------------------------------------------------------------
% GEVORDERD
% ------------------------------------------------------------

\FVDsubsection{Gevorderd}


\begin{oefening}[Evenwicht tussen twee ladingen]{\oefG}

Twee positieve puntladingen:

\[
Q_1=Q
\]

en:

\[
Q_2=4Q
\]

bevinden zich op afstand $d$ van elkaar.

We zoeken op de verbindingslijn een punt waar een positieve
testlading $q$ een resulterende elektrische kracht nul zou
ondervinden.

Noem de afstand van dit punt tot $Q_1$ gelijk aan $x$.

\begin{enumerate}
  \item Leg uit waarom het gezochte punt tussen beide ladingen moet
        liggen.
  \item Teken de twee krachtvectoren op $q$.
  \item Stel een vergelijking op voor $x$.
  \item Los de vergelijking op.
  \item Leg fysisch uit waarom het punt dichter bij de kleinste
        bronlading ligt.
\end{enumerate}

\end{oefening}


\begin{oefening}[Een driehoek van ladingen]{\oefG}

Drie puntladingen bevinden zich op de hoekpunten van een rechthoekige
driehoek.

\[
Q_1=+2{,}0\,\mathrm{\mu C},
\qquad
Q_2=+3{,}0\,\mathrm{\mu C},
\qquad
Q_3=-4{,}0\,\mathrm{\mu C}.
\]

De rechthoekszijden hebben lengtes:

\[
0{,}30\,\mathrm m
\]

en:

\[
0{,}40\,\mathrm m.
\]

Kies zelf een duidelijke plaatsing van de ladingen op een
coördinatenstelsel.

Bepaal vervolgens de resulterende elektrische kracht op $Q_3$.

Je oplossing moet minstens bevatten:

\begin{itemize}
  \item een schets;
  \item de afzonderlijke krachtvectoren;
  \item de gebruikte afstanden;
  \item de componenten;
  \item de grootte van de resulterende kracht;
  \item de richting en zin.
\end{itemize}

\end{oefening}


\begin{oefening}[Ontwerp een configuratie]{\oefG}

Je beschikt over drie puntladingen:

\[
+Q,\qquad +Q,\qquad -q.
\]

Ontwerp een symmetrische configuratie waarbij de resulterende
elektrische kracht op de negatieve lading langs één vooraf gekozen as
ligt.

\begin{enumerate}
  \item Teken je configuratie.
  \item Verklaar welke krachtcomponenten elkaar opheffen.
  \item Verklaar welke componenten elkaar versterken.
  \item Stel een formule op voor de grootte van de resulterende kracht.
\end{enumerate}

\end{oefening}


\begin{oefening}[Modelkritiek]{\oefG}

De wet van Coulomb wordt in dit hoofdstuk gebruikt voor puntladingen.

Een leerling wil dezelfde formule rechtstreeks gebruiken voor twee
grote, onregelmatig gevormde geladen voorwerpen die zeer dicht bij
elkaar staan.

Beoordeel deze aanpak.

Bespreek:

\begin{enumerate}
  \item welke aanname achter het puntladingmodel zit;
  \item waarom die aanname hier problematisch kan zijn;
  \item welke informatie over de ladingsverdeling relevant zou kunnen
        worden;
  \item waarom een model bruikbaar kan zijn zonder universeel geldig te
        zijn.
\end{enumerate}

\end{oefening}


% ============================================================
% 41. ICT-ONDERZOEK
% ============================================================

\FVDsection{Onderzoeken met ICT}

Gebruik een grafische rekenmachine, GeoGebra, Desmos of een
spreadsheet om de functie:

\[
F_e(r)
=
k\frac{|Q_1Q_2|}{r^2}
\]

te onderzoeken.

Neem bijvoorbeeld:

\[
Q_1=2{,}0\,\mathrm{\mu C}
\]

en:

\[
Q_2=3{,}0\,\mathrm{\mu C}.
\]

Onderzoek voor verschillende waarden van $r$:

\begin{enumerate}
  \item hoe de kracht verandert wanneer de afstand verdubbelt;
  \item hoe de kracht verandert wanneer de afstand verdrievoudigt;
  \item bij welke afstand de kracht nog een kwart van haar
        oorspronkelijke waarde bedraagt;
  \item waarom een lineair model niet geschikt is;
  \item hoe je uit de grafiek het inverse-kwadratenkarakter kunt
        herkennen.
\end{enumerate}


Herhaal eventueel voor andere waarden van $Q_1$ en $Q_2$.

Onderzoek welke eigenschappen van de grafiek veranderen en welke
hetzelfde blijven.


% ============================================================
% 42. MISCONCEPTIES
% ============================================================

\FVDsection{Veelvoorkomende denkfouten}

\begin{itemize}

  \item \textbf{``Wrijving maakt elektrische lading.''}

  Nee. Bij gewone elektrostatische processen worden bestaande
  elektrische ladingen gescheiden of verplaatst.

  \item \textbf{``Een positief geladen voorwerp heeft extra protonen
  gekregen.''}

  Meestal niet. In gewone vaste stoffen ontstaat positieve lading
  doorgaans door een tekort aan elektronen.

  \item \textbf{``Een neutraal voorwerp bevat geen lading.''}

  Nee. Het bevat positieve en negatieve ladingen, maar de netto lading
  is nul.

  \item \textbf{``Als twee voorwerpen elkaar aantrekken, moeten ze
  tegengesteld geladen zijn.''}

  Niet noodzakelijk. Een geladen voorwerp kan door ladingsherverdeling
  of polarisatie ook een neutraal voorwerp aantrekken.

  \item \textbf{``Een grotere lading oefent een grotere kracht uit op
  een kleinere lading dan omgekeerd.''}

  Nee. Volgens Newton III zijn de twee interactiekrachten even groot en
  tegengesteld gericht.

  \item \textbf{``Als de afstand verdubbelt, halveert de kracht.''}

  Nee. Volgens de inverse-kwadratenwet wordt de kracht viermaal kleiner.

  \item \textbf{``Een negatief antwoord uit de Coulombformule betekent
  een negatieve krachtgrootte.''}

  Nee. De grootte van een kracht is niet negatief. Gebruik
  $|Q_1Q_2|$ voor de grootte en bepaal richting en zin afzonderlijk.

  \item \textbf{``Bij drie ladingen tel je gewoon alle
  krachtgroottes op.''}

  Nee. Krachten zijn vectoren en moeten vectorieel worden gecombineerd.

\end{itemize}


% ============================================================
% 43. SAMENVATTING
% ============================================================

\FVDsection{Samenvatting}


\FVDsubsection{Elektrische lading}

Elektrische lading wordt voorgesteld door:

\[
q
\qquad\text{of}\qquad
Q.
\]

De SI-eenheid is:

\[
\boxed{
[q]=\mathrm C.
}
\]


Er bestaan positieve en negatieve elektrische ladingen.

Gelijknamige ladingen stoten elkaar af.

Tegengestelde ladingen trekken elkaar aan.


\FVDsubsection{Elementaire lading}

\[
\boxed{
e=1{,}602\times10^{-19}\,\mathrm C.
}
\]

Voor de grootte van een gekwantiseerde lading:

\[
\boxed{
|q|=Ne.
}
\]


\FVDsubsection{Ladingsbehoud}

Voor een elektrisch geïsoleerd systeem:

\[
\boxed{
Q_\mathrm{voor}=Q_\mathrm{na}.
}
\]


\FVDsubsection{Geleiders en isolatoren}

In een geleider kunnen relevante ladingsdragers zich gemakkelijk
verplaatsen.

In een isolator blijven ladingen veel sterker lokaal gebonden.


\FVDsubsection{Influente werking}

Een geladen voorwerp kan zonder contact een herverdeling van lading in
een ander voorwerp veroorzaken.

Dit noemen we elektrostatische influentie.


\FVDsubsection{Wet van Coulomb}

Voor twee puntladingen:

\[
\boxed{
F_e
=
k\frac{|Q_1Q_2|}{r^2}.
}
\]

met:

\[
\boxed{
k\approx8{,}99\times10^9\,
\frac{\mathrm{N\,m^2}}{\mathrm{C^2}}.
}
\]


\FVDsubsection{Afhankelijkheden}

\[
\boxed{
F_e\propto |Q_1Q_2|
}
\]

en:

\[
\boxed{
F_e\propto\frac1{r^2}.
}
\]


\FVDsubsection{Newton III}

\[
\boxed{
\vec F_{1\rightarrow2}
=
-\vec F_{2\rightarrow1}.
}
\]


\FVDsubsection{Superpositie}

Wanneer meerdere ladingen krachten uitoefenen:

\[
\boxed{
\vec F_\mathrm{res}
=
\sum_i\vec F_i.
}
\]


% ============================================================
% 44. WAT MOET JE ZEKER BEHEERSEN?
% ============================================================

\FVDsection{Wat moet je zeker beheersen?}

Om dit onderdeel op het vereiste niveau te beheersen, volstaat het
niet dat je alleen getallen in de wet van Coulomb kunt invullen.

Je moet zeker kunnen:

\begin{itemize}

  \item positieve, negatieve en neutrale lading microscopisch
        verklaren;

  \item werken met coulomb, microcoulomb en nanocoulomb;

  \item de elementaire lading gebruiken;

  \item het behoud van elektrische lading toepassen;

  \item verklaren wat er gebeurt bij laden door wrijving en contact;

  \item het verschil tussen geleiders en isolatoren verklaren vanuit
        de beweeglijkheid van elektrische ladingen;

  \item elektrostatische influentie verklaren;

  \item uitleggen waarom een geladen voorwerp een neutraal voorwerp kan
        aantrekken;

  \item aantrekking en afstoting correct bepalen;

  \item de wet van Coulomb toepassen:

        \[
        F_e=k\frac{|Q_1Q_2|}{r^2};
        \]

  \item elektrische krachtvectoren correct tekenen;

  \item Newton III toepassen op elektrische interacties;

  \item analyseren hoe de kracht verandert wanneer $Q_1$, $Q_2$ of
        $r$ verandert;

  \item het inverse-kwadratenverband herkennen in een formule, tabel of
        grafiek;

  \item een onbekende lading of afstand uit de wet van Coulomb bepalen;

  \item in situaties met meerdere ladingen de afzonderlijke
        krachtvectoren bepalen en vectorieel combineren;

  \item overeenkomsten en verschillen tussen gravitatie en elektrische
        kracht uitleggen;

  \item je resultaat controleren op eenheid, grootteorde, richting en
        fysische betekenis.

\end{itemize}

De oefeningen met de essentieel-markering vormen samen de minimale
oefenset die je moet beheersen.

Let daarbij op: sommige essentiële oefeningen behoren tot
\oefU. Dat is bewust zo. Het vereiste beheersingsniveau vraagt meer dan
alleen standaardberekeningen.


% ============================================================
% 45. VOORUITBLIK
% ============================================================

\FVDsection{Vooruitblik: van kracht naar veld}

We hebben nu twee fundamentele wisselwerkingen onderzocht.

Bij gravitatie:

\[
F_g
=
G\frac{Mm}{r^2}.
\]

Bij elektriciteit:

\[
F_e
=
k\frac{|Qq|}{r^2}.
\]

Bij gravitatie hebben we nog een belangrijke stap gezet.

We beschreven niet alleen de kracht tussen twee massa's.

We zeiden dat een massa $M$ de ruimte rondom zich beïnvloedt door een
gravitatieveld:

\[
M
\longrightarrow
\vec g
\longrightarrow
\vec F_g.
\]

Daarbij:

\[
\vec F_g=m\vec g.
\]

Dat roept een voor de hand liggende vraag op.

Kunnen we hetzelfde doen bij elektrische lading?

Kunnen we de ruimte rond een bronlading $Q$ beschrijven zonder eerst
een tweede lading te kiezen?

Het antwoord is ja.

We zullen een nieuwe vectoriële grootheid invoeren:

\[
\boxed{\vec E}
\]

het \textbf{elektrische veld}.

Dan ontstaat de parallel:

\[
\boxed{
\begin{array}{ccc}
M & \longrightarrow & \vec g
\\
Q & \longrightarrow & \vec E
\end{array}
}
\]

en:

\[
\boxed{
\begin{array}{ccc}
\vec F_g &=&m\vec g
\\[1mm]
\vec F_e&=&q\vec E.
\end{array}
}
\]

In het volgende hoofdstuk onderzoeken we hoe zo'n elektrisch veld eruit
ziet, hoe verschillende ladingen samen een veld vormen en hoe we met
veldlijnen de onzichtbare elektrische invloed in de ruimte zichtbaar
kunnen maken.

\end{document}